\suggestions

Como aporte al continuo proceso de mejora y evolución que caracteriza a cualquier sistema de software,
se sugieren las siguientes acciones:

\begin{itemize}
    \item Desarrollar un \textit{frontend} que consuma el servicio \textit{backend} obtenido 
    como resultado, permitiendo a las instituciones académicas interactuar con los \ac{FM} 
    mediante una interfaz gráfica intuitiva y así captar la variabilidad de sus planes de formación.

    \item Incorporar técnicas de inteligencia artificial y aprendizaje automático 
    para asistir a los gestores curriculares en el proceso de modelado de la 
    variabilidad, sugiriendo estructuras de características, relaciones y 
    restricciones a partir del análisis de planes de estudio existentes o de 
    estándares curriculares de referencia, reduciendo así la curva de aprendizaje 
    asociada al uso de los \ac{FM}.

    \item Incorporar mecanismos de colaboración en tiempo real sobre los \ac{FM}, 
    permitiendo que múltiples gestores académicos trabajen 
    simultáneamente en el diseño curricular, con funcionalidades de bloqueo 
    de elementos bajo edición, fusión automática de cambios y registro detallado 
    de contribuciones para garantizar la trazabilidad y la integridad del modelo.
\end{itemize}
